\chapter{Linux Permission Controls}
\label{ch:5.2}
\chaptertopic{Topic 5.2}

\chapterquote{Every program and every privileged user of the system should operate using the least amount of privilege necessary to complete the job.}{Jerome H. Saltzer and Michael D. Schroeder, ``The Protection of Information in Computer Systems'' (1975)}

\section{Motivation}

Chapter~\ref{ch:5.1} showed what happens when an application trusts input it should not. This
chapter is about a different trust question: once a request is authenticated, how does a system
decide what that identity is actually allowed to do? Every application, database, and operating
system answers this with some access-control model. The three canonical models -- owner-
controlled, centrally labeled, and role-based -- recur constantly in real systems, often blended
together in the same deployment. The chapter then goes one level deeper into the model that
dominates enterprise practice. It closes with what happens when a stolen credential meets a
poorly designed access-control scheme: lateral movement across an entire network.

\begin{apalignment}{Unit 5: Securing Applications and Data}{Topic 5.2 -- Protecting Applications and Data: Managerial Controls and Access Controls}
\begin{apobjectives}
\apobjective{5.2.A}{Mitigate Risk}{Explain how two properties of a data item -- the sensitivity band an organization has assigned it, and the state it happens to be in at the moment, whether resting on a disk, moving across a link, or open in a running process -- together decide which protections apply to it and how strong those protections have to be.}
\apobjective{5.2.B}{Mitigate Risk}{Identify the controls that operate through policy and process rather than through code: classification schemes, named data owners, least privilege and separation of duties, periodic access review, acceptable-use and retention rules, and the vetting an organization applies to the software it buys.}
\apobjective{5.2.C}{Mitigate Risk}{Decide which access control model -- discretionary, mandatory, role-based, or attribute-based -- suits a given body of applications and data, given who must decide on access, how often the answer changes, and how much administrative effort the organization can sustain.}
\apobjective{5.2.D}{Mitigate Risk}{Put an access decision into effect on a Linux system: set ownership and group, choose a permission mode in octal or symbolic form, apply the special bits where they are warranted, control what new files inherit, and reach for a POSIX access control list when the three-class model cannot express the requirement.}
\end{apobjectives}
\end{apalignment}

\section{Data States and Classification}

Two questions decide how much protection a piece of data deserves, and neither is technical. The
first is \emph{what} the data is -- a marketing brochure and a payroll register are both, mechanically,
just bytes on the same disk, but only one creates an obligation if it leaves the building. The
second is \emph{where} the data is at the moment you want to protect it, because the available
mechanisms differ completely by answer.

\begin{definition}{The three states of data}
\textbf{Data at rest} is written to persistent storage -- a database file, a bucket object, a backup
tape -- and is protected by filesystem/object permissions and by storage encryption. \textbf{Data in
transit} is crossing a network link and is protected by transport encryption: TLS for application
traffic, IPsec or SSH tunnels host-to-host, including internal links. \textbf{Data in use} is loaded
into memory and actively being processed, and generally must exist as plaintext for computation to
occur. It is the hardest state to protect, addressed only by process isolation, memory zeroing after
use, hardware enclaves, and, for narrow cases, homomorphic encryption.
\end{definition}

Full-disk encryption is a clarifying example of why the states matter separately. It defends a
powered-down, stolen laptop. But on a running server with the volume already unlocked, every
process that can read the file reads plaintext, and the encryption contributes nothing at that
point -- access control and encryption are complementary controls, not substitutes for each other.
Data in use is precisely the state the Heartbleed vulnerability (CVE-2014-0160, a 2014 buffer
over-read flaw in OpenSSL) attacked.\srccite{5.2}{1} The private keys it disclosed were properly
encrypted at rest and properly protected in transit. They were read directly out of process memory
in between.

A \textbf{classification scheme} assigns each data item to a sensitivity band so handling rules can
be written once per band instead of once per file. Government practice uses Unclassified,
Confidential, Secret, and Top Secret; commercial practice typically uses Public, Internal,
Confidential, and Restricted. The label is administrative, not technical -- nothing about the bytes
of a Restricted file differs from a Public one. But it is what a policy attaches requirements to:
this band may be stored only on encrypted volumes, may not leave a jurisdiction, must be retained
seven years and then destroyed. Classification and state combine rather than substitute. A
Restricted record at rest needs encrypted storage and a tightly scoped permission set. The same
record in transit needs TLS with a modern cipher suite and enforced certificate validation. The same
record in use needs process isolation and a guarantee the plaintext is never written to a log or
crash dump. \textbf{Aggregation} means a collection can warrant a higher classification than any
single member -- one employee's extension is Internal, the complete company directory is a
social-engineering asset. The \textbf{data lifecycle} means classification is not permanent: an
unreleased earnings figure is Restricted until filed, Public afterward. Data no longer needed
should be destroyed on a retention schedule, since data that does not exist cannot be breached.

Controls are conventionally sorted by how they operate. \textbf{Technical controls} are enforced by
machinery -- a permission bit, a firewall rule, an encryption routine. \textbf{Physical controls}
act on the material world -- a locked rack, a badge reader, a shredder. \textbf{Managerial
(administrative) controls} operate through policy, procedure, and human decision. They are what make
the technical controls coherent: a permission bit answers ``may this account read this file,'' but
only a managerial control answers ``should this account exist at all, and who decided so.''

The managerial controls that bear directly on this chapter are: a \textbf{data classification policy}
defining the bands; \textbf{data ownership}, distinguishing the \emph{owner} who decides access, the
\emph{custodian} who operates storage on the owner's behalf, and the \emph{user} who works under the
rules; and \textbf{least privilege} and \textbf{separation of duties}. Also directly relevant: periodic
\textbf{access review} that catches privilege accumulation before it is exploited;
\textbf{onboarding/offboarding} procedures governing when access is granted and revoked;
\textbf{change management} ensuring a permission change is requested, approved, and recorded;
\textbf{training and acceptable-use policy}; and \textbf{vendor/third-party assessment} extending the
same expectations to software the organization did not write itself. Most well-known data losses
were not defeated cryptography. They were an account that should have been closed, a grant nobody
re-examined, or a production database copy that policy never forbade a developer from taking home.

\section{Access Control Models}

\begin{definition}{Access control}
The policy and technical machinery that decides which \textbf{subjects} (users, processes,
services) can perform which \textbf{permissions} (read, write, delete, execute) on which
\textbf{objects} (files, database tables, API endpoints).
\end{definition}

\textbf{Discretionary Access Control (DAC)} lets the owner of an object decide who else may access
it -- ``discretionary'' meaning the decision sits entirely with that owner. Traditional Unix file
permissions are the textbook DAC system: whoever owns a file sets its read/write/execute bits for
owner, group, and others, with no central authority approving the grant. DAC's strength is
flexibility -- an owner can share access instantly without routing through an administrator -- and
its weakness is exactly the same fact viewed differently. Owners can carelessly over-grant, and
nothing enforces an organization-wide minimum standard, so malware running as a user simply
inherits and abuses whatever that user was permitted to touch. DAC also cannot enforce a rule about
\emph{where} data may go once it is read -- nothing stops an authorized reader from copying a
file's contents into a new file they own and make world-readable.

\textbf{Mandatory Access Control (MAC)} removes discretion from the owner entirely. A central
policy authority assigns \textbf{labels} to both subjects and objects -- UNCLASSIFIED, SECRET, TOP
SECRET -- and a reference monitor compares the two on every access. MAC originated in government and
military systems. Its classical formulation, the \textbf{Bell--LaPadula} model, gives two rules
worth stating carefully because they are frequently reversed. \emph{No read up} means a subject may
not read an object above its clearance. \emph{No write down} means a subject may not write to an
object below its clearance -- preventing a cleared user, or malware acting as that user, from
copying Secret content into an Unclassified file. \begin{figure}[h]
\centering
\begin{tikzpicture}[
  lvl/.style={draw=kickstartermaroon, thick, minimum width=6cm, minimum height=0.8cm, align=center, font=\scriptsize},
]
\node[lvl, fill=kickstartermaroon!25] (ts) at (0,1.6) {TOP SECRET};
\node[lvl, fill=kickstartermaroon!12] (s) at (0,0.8) {SECRET};
\node[lvl, fill=kickstartercream] (u) at (0,0) {UNCLASSIFIED};
\node[draw=kickstartermaroon, circle, thick, minimum size=0.5cm, font=\scriptsize] (subj) at (-3.2,0.8) {S};
\draw[-{Latex[length=2mm]}, gray, dashed] (subj) to[bend left=15] (ts.west) node[midway, above, font=\tiny, gray] {no read up};
\draw[-{Latex[length=2mm]}, kickstartermaroon, thick] (subj) to[bend right=10] (u.west) node[midway, below, font=\tiny] {no write down};
\draw[{Latex[length=2mm]}-, kickstartermaroon, thick] (subj) -- (s.west);
\node[font=\tiny, below=0.05cm of subj] {cleared: SECRET};
\end{tikzpicture}
\caption{Bell--LaPadula's two rules for a subject cleared at SECRET: reading TOP SECRET is
forbidden (no read up), and writing down into an UNCLASSIFIED object is forbidden (no write down)
-- the second rule blocks a cleared user, or malware running as one, from laundering sensitive
content into a less-protected file.}
\label{fig:5.2-bell-lapadula}
\end{figure}

SELinux and AppArmor are MAC's Linux-world descendants,
layered on top of the discretionary permission bits. That is why an SELinux denial can block an
operation that \texttt{ls -l} says is permitted: both layers must allow it. MAC's strength is
enforcement no individual user can override; its weakness is administrative overhead, since labeling
every object correctly requires careful policy design up front.

\begin{definition}{Role-Based Access Control (RBAC)}
An access model where permissions attach to named \textbf{roles} rather than to individual users or
sensitivity labels, and users inherit whatever permissions their assigned role or roles carry.
\end{definition}

RBAC is today's dominant enterprise model precisely because it scales: a hospital with five
thousand staff might define thirty roles -- Nurse, Billing Clerk, Auditor -- rather than five
thousand separate permission lists. When a nurse leaves and a replacement joins, an
administrator reassigns one role rather than reconstructing an individual grant from scratch. RBAC
also enforces separation of duties naturally, since a data-entry role can be kept structurally
distinct from a data-approval role. Its failure mode is \textbf{role proliferation}: an
organization that creates ever more fine-grained roles (``Junior Billing Clerk in Region~4'')
recreates DAC's original complexity problem in a different shape.

\begin{definition}{Attribute-Based Access Control (ABAC)}
An access model that evaluates a policy expression at request time over attributes of the subject,
the object, the action, and the environment, rather than consulting a fixed list of grants or role
assignments.
\end{definition}

ABAC can express conditions the other three models cannot. It can permit a clinician to read a
record only if four conditions all hold: the clinician is assigned to the patient's care team; the
record's classification is at or below the clinician's authorization; the request originates from a
managed device on the hospital network; and the current time falls within the clinician's shift.
All four conditions are evaluated live, without inventing a distinct role for every combination --
the \textbf{role explosion} that afflicts large RBAC deployments. The cost is that policies become software: they need authoring, testing, and
version control, the attribute sources must stay authoritative and current, and a wrong attribute
silently produces a wrong decision. NIST SP 800-162 is the standard reference for ABAC.\srccite{5.2}{2}

\begin{lstlisting}[language=Python, caption={Enforcing DAC, MAC, RBAC, and ABAC access decisions}]
LABEL_RANK = {"unclassified": 0, "secret": 1, "top_secret": 2}

def dac_allowed(owner, requester, acl):
    """acl: set of user names the owner has explicitly granted access to."""
    return requester == owner or requester in acl

def mac_allowed(subject_label, object_label):
    return LABEL_RANK[subject_label] >= LABEL_RANK[object_label]

def rbac_allowed(user_roles, role_permissions, requested_permission):
    """user_roles: set of role names. role_permissions: dict role -> set of
    permission strings."""
    granted = set()
    for role in user_roles:
        granted |= role_permissions.get(role, set())
    return requested_permission in granted

def abac_allowed(subject, obj, action, environment, policy):
    """policy: a callable(subject, obj, action, environment) -> bool,
    evaluated live against current attributes rather than a stored grant."""
    return policy(subject, obj, action, environment)
\end{lstlisting}

In practice the models compose rather than compete. A single enterprise commonly runs RBAC for
application entitlements, DAC on the filesystems those applications use, MAC through SELinux to
confine daemons, and ABAC at the API gateway for context-dependent conditions. The design question
is not which model is universally best but which decision each layer is being asked to make.

In practice, systems frequently blend all three: an operating system enforces MAC-style clearance
labels overall while an application built on top of it manages job functions through RBAC, with DAC
as the fallback for user-owned files. Cutting across every model is \textbf{least privilege}: a web
server process reading a public directory should hold no write or delete permission at all, and an
intern account should never reach salary tables. Least privilege is what caps the
\textbf{blast radius} once any single account is compromised -- an attacker who hijacks a
low-privilege account inherits only that account's narrow reach, not the whole system.

\begin{misconception}
\textbf{``DAC is weak, so only MAC is secure.''} Each model fits a different context. MAC's
overhead is impractical for ordinary consumer software; DAC is entirely appropriate when owners can
be trusted to manage their own data reasonably.
\end{misconception}

\section{Configuring Access Control on Linux}

Every file and directory on a Linux system carries an \textbf{inode} recording its owner, its
group, and a compact permission field. The model divides potential accessors into exactly three
classes -- \textbf{owner}, \textbf{group}, and \textbf{other} (everyone else, including, on an
internet-facing server, an unauthenticated web-server worker process). It tracks three
independent bits per class: \textbf{read} (4), \textbf{write} (2), and \textbf{execute} (1). Running
\texttt{ls -l} displays this as a ten-character string, e.g.\ \texttt{-rw-r-----}. The leading
character is the file type (\texttt{-} regular file, \texttt{d} directory, \texttt{l} symlink). Each
of the remaining three triads belongs to owner, group, and other in that order -- a detail
students routinely miscount by including the type character as part of the owner's bits.

\begin{lstlisting}[caption={Reading and setting Linux permissions}]
$ ls -l /srv/app/config.yml
-rw-r----- 1 appsvc appadm 2048 Sep 2 09:14 /srv/app/config.yml
# owner appsvc: read+write; group appadm: read only; other: nothing -> mode 640

$ chmod 755 /usr/local/bin/deploy.sh   # rwxr-xr-x: typical executable script
$ chmod 640 /srv/app/config.yml        # rw-r-----: typical config/data file
$ chmod u+x script.sh                  # symbolic: add execute for owner only
$ chmod g-w,o= shared.txt              # remove group write; clear other entirely
\end{lstlisting}

Two evaluation rules are easy to get wrong. First, the three triads are checked in order and the
\emph{first match wins}: if the requesting process is the file's owner, only the owner triad is
consulted. So a mode of \texttt{077} denies the owner everything while granting everyone else full
access. Second, on a \emph{directory} the bits mean something different from a file: \texttt{r}
permits listing names inside it, \texttt{w} permits creating and deleting entries, and \texttt{x}
(the search bit) permits traversing into it to reach a named child even without \texttt{r}.
Because \texttt{w} on a directory governs the directory's entry list rather than the files' contents,
holding write permission on a directory is enough to delete a file inside it that you cannot read
and do not own. That is the exact problem the \textbf{sticky bit} (\texttt{chmod +t}, seen as \texttt{t}
in \texttt{/tmp}'s mode) exists to solve, by restricting deletion inside a shared, world-writable
directory to a file's own owner. The related \textbf{setuid} and \textbf{setgid} bits make a
program run with the privileges of its owning user or group rather than the invoking user's. That
makes them powerful, and consequently a classic privilege-escalation target when set on a file an
attacker can also modify.

Permissions on a newly created file are not chosen by the creating program alone. The kernel
computes the mode as the process's requested mode (conventionally 666 for a file, 777 for a
directory) with every bit set in the process's \textbf{umask} cleared.

\begin{lstlisting}[caption={umask determines what new files inherit}]
$ umask 0022
$ touch newfile && ls -l newfile | cut -d' ' -f1
-rw-r--r--          # 666 with 022 cleared

$ umask 077
$ touch private && ls -l private | cut -d' ' -f1
-rw-------          # nothing for group or other
\end{lstlisting}

A umask is a default, not an enforcement mechanism -- it decides what a file gets at birth, and any
process may \texttt{chmod} afterward. It is also per-process and inherited, so setting it in an
interactive shell has no effect on a service started by the system's service manager. That service
needs its own mask configured directly (e.g.\ \texttt{UMask=} in a systemd unit file).

The three-class model is intentionally coarse: it cannot express ``grant write to exactly these two
users but not the rest of the group,'' nor different default permissions for different creators
within one directory. A \textbf{POSIX Access Control List (ACL)}, managed with \texttt{getfacl} and
\texttt{setfacl}, layers arbitrarily many additional user- and group-specific grants on top of the
base owner/group/other triad without replacing it. ACLs are additive, so the base bits are always
evaluated too. A trailing \texttt{+} on an \texttt{ls -l} mode string signals that an ACL is present
and the nine bits are not the whole story; auditing permissions without noticing it produces a
confidently wrong conclusion.

\begin{misconception}
\textbf{``If \texttt{ls -l} shows the permission, the operation is allowed.''} On a system running
SELinux or AppArmor, a MAC policy can still deny an operation the discretionary bits permit -- both
layers must agree, and a trailing \texttt{.} on the mode string signals an attached security context
worth checking separately.
\end{misconception}

\section{Roles, Permissions, and Least Privilege}

A \textbf{permission} is a single narrow grant -- ``read the invoices table,'' ``restart the web
service.'' Real systems have hundreds or thousands of possible permissions, and assigning them one
at a time to every individual user does not scale. It is slow, error-prone, and nearly impossible
to audit after the fact (``does Dana still need delete-access to payroll from a project two years
ago?''). RBAC's roles are the answer -- bundling permissions into named units a user is assigned as
a whole, so a job change means reassigning a role rather than hand-editing a scattered permission
list.

\begin{example}[Why over-granting is silent and under-granting is loud]
If a new employee's role is missing a permission they actually need, they notice immediately and
complain -- the gap is loud and gets fixed fast. If a role carries a permission the employee will
never use, nothing breaks and nobody notices -- the excess is silent. This asymmetry is exactly why
\textbf{privilege creep} accumulates unopposed. Over months and years, employees who change teams or
take on temporary projects gain more and more permissions, because access is added when needed but
almost never removed once it stops being needed. Removal requires someone to notice a gap
that, by construction, never announces itself.
\end{example}

The security stake here is not tidiness -- it is blast radius. A support agent's role scoped tightly
to ticket-handling permissions means a phished credential lets an attacker read and reply to
tickets: contained, if unwelcome. The same account, after years of unchecked privilege creep,
carrying \texttt{delete\_invoice} and \texttt{read\_all\_customer\_pii} because a manager once said
``just give them access to be safe,'' turns the identical phishing email into a much larger breach.
\textbf{Role review (recertification)} is the deliberate counterweight: on a periodic schedule,
managers or system owners must explicitly re-confirm that every person assigned to a role still
needs it, and anything not re-confirmed is revoked.

\begin{lstlisting}[language=Python, caption={Flagging users whose role permissions exceed their job requirement}]
def find_overprivileged_users(users, role_permissions, job_requirements):
    """users: dict user -> set of assigned role names.
       role_permissions: dict role -> set of permission strings.
       job_requirements: dict user -> set of permissions actually needed.
       Returns sorted usernames holding a permission beyond their
       job requirement."""
    flagged = []
    for user, roles in users.items():
        granted = set()
        for role in roles:
            granted |= role_permissions.get(role, set())
        excess = granted - job_requirements.get(user, set())
        if excess:
            flagged.append(user)
    return sorted(flagged)
\end{lstlisting}

\textbf{Separation of duties} reinforces least privilege structurally: the person who \emph{creates}
a vendor payment should never also be the person who \emph{approves} it. That way, completing a fraud
requires collusion between two people rather than the abuse of a single over-broad role. The most
common design mistake here is a single giant ``power user'' role bundling every permission a manager
might conceivably need across every function, rather than composing narrower roles (Approver,
Report Viewer, Team Editor) assigned only as actually needed. It takes more setup effort up front,
but it is the difference between answering ``who can delete invoices?'' by listing one role's
members versus manually auditing a thousand individual accounts.

\begin{misconception}
\textbf{``Giving extra access `just in case' is harmless if it's never used.''} Unused excess
permission is exactly what increases blast radius the moment an account is compromised. It is a
silent risk, not a neutral one, regardless of how rarely it is actually exercised.
\end{misconception}

\section{Credential Theft and Lateral Movement}

Most large breaches do not end where they begin. A phished marketing employee's laptop usually has
no direct access to a finance database or a domain controller -- what makes a breach catastrophic is
what happens next.

\begin{definition}{Lateral movement}
An attacker using a compromised account or host to reach additional systems within a network after
an initial foothold, typically by harvesting further credentials at each stop.
\end{definition}

Credential theft is lateral movement's starting fuel: once an attacker has one foothold, they
harvest passwords, password hashes, or authentication tokens sitting in memory, configuration
files, or the OS credential cache. Crucially, an attacker frequently does not need the plaintext
password at all. Many enterprise authentication protocols accept \textbf{hash-based} or
\textbf{ticket-based} authentication, so a stolen hash or session ticket can be replayed to
authenticate as that user elsewhere without ever cracking it. This attack family is known as
``pass-the-hash'' or ``pass-the-ticket.''

\begin{example}[How least privilege stops lateral movement in its tracks]
A stolen credential belonging to a normal user with narrow permissions stalls quickly -- this is
least privilege from earlier in this chapter doing exactly its job. If the compromised account is
an IT administrator, or if the organization reuses one local-administrator password across every
machine (a common and genuinely dangerous shortcut), the attacker instead hops from machine to
machine nearly unopposed. The attacker harvests more credentials at each stop, hunting for a path
to domain controllers, backup servers, or financial databases. This combination of lateral movement and
\textbf{privilege escalation} expands an attacker's reach both sideways and upward at once.
\end{example}

Several access-control practices from this chapter directly blunt this chain. \textbf{Multi-factor
authentication} breaks the ``stolen password equals access'' assumption outright, since even a
perfectly stolen credential is useless without the second factor, stopping the vast majority of
credential-replay attempts cold. \textbf{Credential hygiene} -- unique passwords per system, no
shared local-admin password across machines, no plaintext credentials in scripts -- removes the
``one stolen credential unlocks everything'' shortcut that fuels rapid lateral movement in the first
place. \textbf{Least privilege and network segmentation} limit how far a single stolen credential
can physically travel, exactly as Chapter~\ref{ch:3.3} described.

\begin{lstlisting}[language=Python, caption={Detecting a lateral-movement pattern in an authentication log}]
def flag_lateral_movement(events, host_threshold=5, window_seconds=600):
    """events: list of dicts like {"account": str, "dst_host": str,
       "timestamp": int}. Returns sorted accounts that authenticated
       to >= host_threshold distinct hosts within any window_seconds
       window."""
    from collections import defaultdict
    by_account = defaultdict(list)
    for e in events:
        by_account[e["account"]].append((e["timestamp"], e["dst_host"]))

    flagged = set()
    for account, entries in by_account.items():
        entries.sort()
        left = 0
        hosts_in_window = []
        for right in range(len(entries)):
            hosts_in_window.append(entries[right])
            while entries[right][0] - entries[left][0] > window_seconds:
                left += 1
            distinct_hosts = {h for _, h in entries[left:right + 1]}
            if len(distinct_hosts) >= host_threshold:
                flagged.add(account)
                break
    return sorted(flagged)
\end{lstlisting}

Detection matters as much as prevention, because a sufficiently patient attacker will eventually
find some foothold. Security teams watch for the same behavioral-anomaly signatures introduced in
Chapter~\ref{ch:3.5} and Chapter~\ref{ch:4.4}. Three signatures matter most: a single account
authenticating to an unusually large number of distinct hosts in a short window; authentication at
odd hours relative to the user's normal pattern; or a user account suddenly invoking administrative
tools it has never touched before.

\begin{misconception}
\textbf{``A strong password alone stops credential theft.''} Many lateral-movement techniques
replay a stolen hash or ticket without ever needing the plaintext password at all. MFA and
credential hygiene matter more here than password strength alone.
\end{misconception}

\section*{Review Questions}
\begin{reviewquestions}
\item A government system enforces that a user cleared to SECRET cannot read a TOP SECRET document
even if the document's owner wants to share it. Identify which access control model this is and
explain why DAC could not enforce this same guarantee.
\item An RBAC system has a single ``Manager'' role bundling every permission a manager in any
department might conceivably need. Explain why this violates least privilege even though every
individual manager is, in fact, a manager.
\item Explain why privilege creep is described as an asymmetric problem, using the loud-gap versus
silent-excess distinction from this chapter.
\item An attacker steals a password hash from a compromised laptop and uses it to authenticate to
three other machines without ever learning the plaintext password. Explain why MFA would have
stopped this specific attack while a stronger password policy alone would not have.
\item Explain how least privilege and network segmentation work together to contain lateral
movement, using the marketing-employee-versus-domain-admin scenario from this chapter.
\item A file shows mode \texttt{-rw-r-----} and its owning directory is world-writable. Explain,
using the directory-write rule from this chapter, how another user could delete this file despite
having no permission bits on the file itself, and name the mechanism that exists specifically to
prevent this in a shared directory.
\item A hospital wants clinicians to read a patient record only if they are on that patient's care
team, on a managed device, and during their shift. Explain why RBAC cannot express this condition
without role explosion, and name the model that can.
\item Using the Bell--LaPadula diagram, explain why ``no write down'' is necessary in addition to
``no read up'' -- specifically, describe an attack a system enforcing only no-read-up would remain
vulnerable to.
\item Malware executing under a SECRET-cleared user's session attempts to copy a TOP SECRET file's
contents into a new UNCLASSIFIED file the malware creates. Explain which Bell--LaPadula rule blocks
this, and why the block applies even though the malware is running with a legitimately cleared
user's own privileges rather than a stolen or escalated identity.
\item An unreleased quarterly earnings figure is classified Restricted. Explain what protections it
needs while at rest on a finance server, while in transit to the company's outside auditors, and
while open in a spreadsheet on an analyst's laptop, and explain why the classification band alone
does not tell you the answer to any of the three.
\end{reviewquestions}

\begin{answerkey}
\item This is Mandatory Access Control, applying Bell--LaPadula's no-read-up rule. DAC could not
enforce it because DAC leaves the access decision to the object's owner, who could simply choose to
share the document regardless of the requester's clearance. MAC works precisely because the policy
is set centrally and cannot be overridden by any individual owner's discretion.
\answerlabel{5.2.C}
\item Least privilege requires that an account hold only the access its actual function demands.
Bundling every permission any manager anywhere might conceivably need into one role grants each
individual manager far more access than their specific job requires. So a phished or malicious
manager account has a blast radius sized to the whole bundled role rather than to that manager's
actual duties.
\answerlabel{5.2.B}
\item A missing permission a new employee needs is loud: they notice immediately and report it, so
it gets fixed quickly. An unused excess permission is silent: nothing breaks and nobody notices, so
it persists indefinitely. Because access is added on demand but almost never removed without a
deliberate review process, this asymmetry causes privilege to accumulate one-directionally over
time.
\answerlabel{5.2.B}
\item Pass-the-hash authentication replays the stolen hash directly without ever needing the
plaintext password, so a stronger password policy changes nothing about how crackable the plaintext
is -- the attacker never attempts to crack it. MFA stops the attack because the stolen hash alone is
insufficient to complete authentication without the second factor, which the attacker does not
possess.
\answerlabel{5.2.D}
\item Least privilege ensures a compromised marketing account's credentials only unlock what a
marketing job actually requires. Network segmentation ensures that even if the attacker tries
to reach further, the finance database and domain controller sit on network segments the
marketing host cannot route to in the first place. One control limits what the account is
permitted to do, the other limits what hosts it can even reach, and together they contain lateral
movement in two independent dimensions.
\answerlabel{5.2.B}
\item Because the directory carries write permission for other, any user can create or delete
entries in its listing regardless of the permissions on the individual files inside it. Deleting a
file only requires write permission on its parent directory, not on the file itself. The sticky bit,
set with \texttt{chmod +t}, exists specifically to prevent this in a shared directory by restricting
deletion to a file's own owner (or the directory's owner, or root) even when the directory itself is
world-writable. That is exactly why \texttt{/tmp} carries it by default.
\answerlabel{5.2.D}
\item RBAC's roles are static grants: a role either has a permission or it does not, independent of
time, location, device, or the specific record in question. Expressing ``only during this shift, only
on a managed device, only for a patient on this clinician's care team'' would require a distinct role
for every combination of those conditions -- role explosion. Attribute-Based Access Control (ABAC)
evaluates such conditions live, as a policy expression over the request's actual attributes, so it
can express this rule directly without enumerating roles for every combination.
\answerlabel{5.2.C}
\item A system enforcing only no-read-up still allows a subject cleared at SECRET to write freely
into any object, including one labeled UNCLASSIFIED. A legitimately cleared user, or malware
running under that user's session, could read a SECRET or TOP SECRET file -- permitted, since
reading at or below one's own clearance is always allowed. That user or malware could then write
the contents into a new UNCLASSIFIED file. That effectively launders sensitive content down to a
label anyone can read, with no rule in place to stop the write itself. No-write-down closes exactly
this gap by forbidding a subject from writing to any object below its own clearance level.
\answerlabel{5.2.C}
\item No-write-down blocks it: the malware is writing TOP SECRET content into an UNCLASSIFIED
object, which is below the acting subject's SECRET clearance regardless of what that content
actually is. The block applies even though the malware runs under a legitimately cleared user's
own privileges, because Bell--LaPadula's reference monitor evaluates every access by comparing
labels. It does not evaluate whether the specific action is intentional or malicious. MAC's entire
strength is that it constrains what an account is permitted to do structurally, so it protects
against a legitimate account behaving maliciously (through malware or otherwise) exactly as well as
it protects against a mismatched external request.
\answerlabel{5.2.C}
\item At rest, a Restricted item needs encrypted storage and a tightly scoped permission set that
restricts read access to authorized users. In transit, it needs TLS with a modern cipher suite and
enforced certificate validation. In use, it needs process isolation and a guarantee the plaintext is
never written to a log or crash dump. The classification band alone only says how sensitive the data
is; it says nothing about which specific technical mechanism applies, because that mechanism is
determined by the item's current state. Sensitivity band and current state have to be combined to
decide both which protections apply and how strong they have to be.
\answerlabel{5.2.A}
\end{answerkey}

\section*{Sources}
\begin{enumerate}[label=\lbrack\arabic*\rbrack, leftmargin=2.2em]
\item \srcitem{5.2}{1}National Vulnerability Database, CVE-2014-0160 (``Heartbleed''), U.S. National
Institute of Standards and Technology, published April 7, 2014 -- a buffer over-read in the TLS/DTLS
heartbeat handling of OpenSSL 1.0.1 through 1.0.1f, allowing an attacker to read private keys and
other sensitive data directly out of process memory.
\item \srcitem{5.2}{2}Vincent C. Hu, David Ferraiolo, Richard Kuhn, et al., \emph{Guide to Attribute
Based Access Control (ABAC) Definition and Considerations}, NIST Special Publication 800-162,
National Institute of Standards and Technology, originally published January 2014, updated August
2, 2019.
\end{enumerate}
